\section{Studi Kasus: Verifikasi Algoritma Pengurutan}

Algoritma pengurutan (\textit{sorting}) merupakan salah satu topik paling fundamental dalam ilmu komputer. Verifikasi formal algoritma pengurutan menyediakan studi kasus yang lebih kompleks dibandingkan pencarian linear, karena melibatkan loop bersarang (\textit{nested loops}) dan properti yang lebih rumit \cite{rosen2019}.

\subsection{Insertion Sort: Algoritma}

\begin{verbatim}
    InsertionSort(A, n):
        for j := 1 to n-1 do
            key := A[j]
            i := j - 1
            while i >= 0 and A[i] > key do
                A[i+1] := A[i]
                i := i - 1
            A[i+1] := key
\end{verbatim}

\textbf{Spesifikasi}:
\begin{itemize}
    \item \textbf{Prekondisi}: $n \geq 1$ dan $A[0..n-1]$ terdefinisi
    \item \textbf{Postkondisi}: $A[0] \leq A[1] \leq \ldots \leq A[n-1]$ dan $A$ merupakan permutasi dari array masukan
\end{itemize}

\subsection{Loop Invariant untuk Loop Luar}

\begin{definisi}[Loop Invariant Insertion Sort]
Invarian loop luar: ``Pada awal setiap iterasi loop \texttt{for} dengan indeks $j$, subarray $A[0..j-1]$ terurut secara tidak menurun dan merupakan permutasi dari elemen-elemen asli $A[0..j-1]$.''
\end{definisi}

\textbf{Verifikasi:}

\textbf{1. Inisialisasi} ($j = 1$): Subarray $A[0..0]$ hanya berisi satu elemen, yang secara trivial terurut. $\checkmark$

\textbf{2. Pemeliharaan}: Asumsikan $A[0..j-1]$ terurut. Badan loop mengambil $A[j]$ (disebut $key$) dan menyisipkannya ke posisi yang tepat dalam $A[0..j-1]$ melalui loop dalam (\texttt{while}). Elemen-elemen yang lebih besar dari $key$ digeser satu posisi ke kanan. Setelah penyisipan, $A[0..j]$ terurut dan merupakan permutasi dari elemen asli. $\checkmark$

\textbf{3. Terminasi} ($j = n$): Invarian menyatakan $A[0..n-1]$ terurut dan merupakan permutasi dari array masukan. Ini adalah postkondisi. $\checkmark$

\subsection{Loop Invariant untuk Loop Dalam}

Loop dalam (\texttt{while}) juga memiliki invariannya sendiri:

\begin{konsep}
Invarian loop dalam: ``$A[i+2..j]$ berisi elemen-elemen yang lebih besar dari $key$ dalam urutan terurut, dan $key$ akan disisipkan pada posisi $i+1$.''

Fungsi pengurangan loop dalam: $V = i + 1$. Karena $i$ menurun 1 pada setiap iterasi dan loop berhenti ketika $i < 0$ atau $A[i] \leq key$, loop dalam pasti berterminasi.
\end{konsep}

\subsection{Properti yang Harus Dijaga}

Perhatikan bahwa postkondisi pengurutan memiliki dua komponen:
\begin{enumerate}
    \item \textbf{Terurut}: $A[0] \leq A[1] \leq \ldots \leq A[n-1]$.
    \item \textbf{Permutasi}: Array hasil merupakan permutasi dari array masukan (tidak ada elemen yang hilang atau ditambahkan).
\end{enumerate}

Properti kedua sering diabaikan oleh pemula tetapi sangat penting. Program yang mengganti seluruh elemen array menjadi $0$ menghasilkan output terurut, tetapi jelas salah karena bukan permutasi dari masukan.
